%%%%%%%%%%  scan idx 208  |  folio 195  %%%%%%%%%%
\SGAPrintPage{195}%
\sgaStableTarget{sga6.e02.appendix-i-a-counterexample-of-verdier}\section*{APPENDIX I\\[3pt] A COUNTEREXAMPLE OF VERDIER\protect\SGAAuthorFootnote{1}{Apart from a few details of presentation, this appendix reproduces two letters from J.-L. Verdier (30/9/67 and 11/10/67) to the editor, in reply to a question from the latter.}}

\sgaStableTarget{sga6.e02.sec.0}\subsection*{0. Introduction}
For a scheme $S$, we denote by $\Mod(S)$ (resp. $\Qcoh(S)$) the category of $\underline{O}_S$-Modules (resp. quasi-coherent $\underline{O}_S$-Modules). The purpose of this appendix is to draw the reader's attention to certain pathological properties of the category $\mathrm D(\Qcoh(S))$, and in particular to prove the following facts:

\sgaStableTarget{sga6.e02.para.0-1}0.1. Let $X$ be an affine scheme with ring $A$, and let $I$ be an injective $A$-module. We know ([H] II 7.14) that, when $A$ is noetherian, the sheaf $\widetilde I$ is an injective $\underline{O}_X$-Module. When $A$ is not noetherian, $\widetilde I$ need not be injective, or even flasque, even when the underlying space of $X$ is noetherian and of finite Krull dimension. Moreover, the restriction of $\widetilde I$ to an open subset $U$ of $X$ is not in general an injective object of $\Qcoh(U)$.

\sgaStableTarget{sga6.e02.para.0-2}0.2. Let $f:X\longrightarrow Y$ be a morphism of quasi-compact and quasi-separated schemes. The functor $\underline{\underline{R}}f_*:\mathrm D^+(X)\longrightarrow\mathrm D^+(Y)$ does not in general induce on $\mathrm D^+(\Qcoh(X))$ (a full subcategory of $\mathrm D^+(X)$ by (II 3.5)) the right derived functor of $f_*:\Qcoh(X)\longrightarrow\Mod(Y)$. More precisely, an injective object of $\Qcoh(X)$ need not be acyclic for $f_*:\Mod(X)\longrightarrow\Mod(Y)$.

\sgaStableTarget{sga6.e02.para.0-3}0.3. Let $S$ be a quasi-compact and quasi-separated scheme. An object of $\Qcoh(S)$, even an injective one, need not be acyclic for the coherer (II 3.2). The functor $\phi:\mathrm D^b(\Qcoh(S))\longrightarrow\mathrm D^b(S)$ (II 3.0) is not in general fully faithful.

\sgaStableTarget{sga6.e02.sec.1.occ2}\subsection*{1. The counterexamples}
We first recall the following lemma (SGA 2 II 9):

\sgaStableTarget{sga6.e02.lem.1-1}\textbf{Lemma 1.1.} \uline{Let} $A$ \uline{be a ring,} $X=\Spec(A)$, $\underline{f}=(f_\alpha)$ \uline{a finite system of elements of} $A$, $Y$ \uline{the closed subset defined by} $\underline{f}(x)=0$\uline{. For} $i>0$\uline{, the following conditions are equivalent:}
\begin{enumerate}[label=(\roman*)]
\item \uline{for every injective} $A$\uline{-module} $I$\uline{, we have} $H^i_Y(X,\widetilde I)=0$;
\item \uline{the pro-object} $n\mapsto H_i(\underline{f}^n,A)$ \uline{is zero, i.e. for every} $n$\uline{, there exists} $n'\geq n$
%%%%%%%%%%  scan idx 209  |  folio 196  %%%%%%%%%%
\SGAPrintPage{196}%
\uline{such that}
\[
 H_i(\underline{f}^{n'},A)\longrightarrow H_i(\underline{f}^n,A)
\]
\uline{is zero}.
\end{enumerate}

\sgaStableTarget{sga6.e02.para.1-2}1.2 We shall use (1.1) to prove (0.1). Let $k$ be a field, $B=k[[x,y]]$, $\underline{f}$ the system $(x,y)$, $M=\coprod_{n>0}B/(\underline{f}^n)$, $A$ the $B$-algebra $D_B(M)$ (EGA $0_{\mathrm{IV}}$ 18.2.3). We have trivially
\[
 H_2^A(\underline{f}^n,A)\simeq H_2^B(\underline{f}^n,A).\qquad\text{\SGAAuthorFootMark{1}}
\]
\SGAAuthorFootText{1}{The superscript A or B indicates the ring over which the Koszul complex is computed.}
On the other hand, we have
\[
 H_2^B(\underline{f}^n,A)\simeq H_2^B(\underline{f}^n,B)\oplus H_2^B(\underline{f}^n,M)
\]
but $H_2^B(\underline{f}^n,B)=0$ since $\underline{f}$ is a regular sequence (EGA III 1.1.4 and $0_{\mathrm{IV}}$ 15.1.20), so we have
\[
 H_2^B(\underline{f}^n,A)\simeq H_2^B(\underline{f}^n,M).
\]
Finally, we have (EGA III 1.1.3.3)
\[
 H_2^B(\underline{f}^n,M)\simeq H_B^0(\underline{f}^n,M),
\]
where $H_B^0(\underline{f}^n,M)$ is canonically identified with the submodule ${}_{(\underline{f}^n)}M$ of $M$ annihilated by $(\underline{f}^n)$. Hence finally:
\sgaStableTarget{sga6.e02.eq.unnamed.occ6}\begin{equation}\tag{*}
 H_2^A(\underline{f}^n,A)\simeq {}_{(\underline{f}^n)}M.
\end{equation}
Since the transition morphism $H_2^A(\underline{f}^{n'},A)\longrightarrow H_2^A(\underline{f}^n,A)$ is induced by multiplication by $\underline{f}^{\,n'-n}$ (EGA III 1.1.6), we check immediately, using (*), that \uline{the pro-object} $n\mapsto H_2^A(\underline{f}^n,A)$ \uline{is not zero}. From (1.1) we then deduce that there exists an injective $A$-module $I$ such that $H^2_Y(X,\widetilde I)\ne0$, and thus such that $\widetilde I$ is not flasque (SGA 2 I 2.12). We note in passing that the projection $\Spec(A)\longrightarrow\Spec(B)$ is a homeomorphism,
%%%%%%%%%%  scan idx 210  |  folio 197  %%%%%%%%%%
\SGAPrintPage{197}%
so the underlying space of $X$ is noetherian of dimension $2$. Let $I$ be an injective $A$-module such that $H^2_Y(X,\widetilde I)\ne0$, and denote by $U$ the open complement of $Y$ in $X$. We have (SGA 2 II 4)
\[
 H^2_Y(X,I)\simeq H^1(U,I)\ne0.
\]
Now, since $U$ is separated, the functor $\mathrm D^+(\Qcoh(U))\to\mathrm D^+(U)$ is fully faithful (II 3.5), so $\widetilde I|U$ is not injective in $\Qcoh(U)$. We have therefore proved assertion (0.1).

Denote by $Z$ the quasi-compact and quasi-separated scheme obtained by gluing two copies of $X$ along $U$, and by $j_1,j_2$ the two canonical open immersions of $X$ into $Z$, so that we have a cartesian square
\[
\begin{tikzcd}
 U \arrow[r,"i"] \arrow[d,"i"'] & X \arrow[d,"j_1"]\\
 X \arrow[r,"j_2"] & Z .
\end{tikzcd}
\]
We have
\[
 j_2^*R^1j_{1*}(\widetilde I)\simeq R^1i_*(i^*\widetilde I),
\]
whence
\[
 H^0(X,R^1j_{1*}(\widetilde I))\simeq H^1(U,\widetilde I)\ne0,
\]
so $R^1j_{1*}(\widetilde I)\ne0$, which proves (0.2).

On the other hand, by (II 3.4 and 3.5), we have
\[
 \underline{\underline{R}}Q\,\underline{\underline{R}}j_{1*}(\widetilde I)\simeq j_{1*}(\widetilde I).
\]
We deduce a spectral sequence whose exact sequence of low-degree terms gives the isomorphism
\[
 R^1j_{1*}(\widetilde I)\simeq R^2Q(j_{1*}\widetilde I).
\]
Since the sheaf $j_{1*}\widetilde I$ is injective in $\Qcoh(Z)$, we have proved the first
%%%%%%%%%%  scan idx 211  |  folio 198  %%%%%%%%%%
\SGAPrintPage{198}%
part of (0.3).

Finally, let $E\in\Ob\Qcoh(Z)$. We have the duality isomorphism
\[
 \underline{\underline{R}}\Hom_{\Mod}(E,\underline{\underline{R}}j_{1*}(\widetilde I))
 \simeq \underline{\underline{R}}\Hom_{\Mod}(j_1^*E,\widetilde I),
\]
but, since $X$ is affine, we have (II 3.5)
\[
 \underline{\underline{R}}\Hom_{\Mod}(j_1^*E,\widetilde I)
 \simeq \underline{\underline{R}}\Hom_{\Qcoh}(j_1^*E,I)
 \simeq \Hom(j_1^*E,\widetilde I).
\]
We deduce
\[
 \underline{\underline{R}}\Hom_{\Mod}(E,\underline{\underline{R}}j_{1*}(\widetilde I))
 \simeq \Hom(j_1^*E,\widetilde I),
\]
whence a spectral sequence whose exact sequence of low-degree terms gives the isomorphism
\[
 \Ext^2_{\Mod}(E,j_{1*}\widetilde I)
 \simeq \Hom(E,R^1j_{1*}\widetilde I).
\]
Consequently, the identity morphism of $E=R^1j_{1*}(\widetilde I)$ gives a nonzero element of $\Ext^2_{\Mod}(E,j_{1*}\widetilde I)$, so the functor $\mathrm D^b(\Qcoh(Z))\longrightarrow \mathrm D^b(Z)$ is not fully faithful, which completes the proof of (0.3).

% Original English continuation definitions, retained in their original order.
\providecommand{\Diff}{\operatorname{Diff}}
\providecommand{\Hom}{\operatorname{Hom}}
\providecommand{\im}{\operatorname{Im}}
\providecommand{\ind}{\operatorname{ind}}
\providecommand{\Cone}{\operatorname{Cone}}
\providecommand{\Loclib}{\operatorname{Loclib}}
\providecommand{\ob}{\operatorname{ob}}
\providecommand{\Parf}{\operatorname{Parf}}
\providecommand{\Mod}{\operatorname{Mod}}
\providecommand{\Ob}{\operatorname{Ob}}
\providecommand{\Coker}{\operatorname{Coker}}
\providecommand{\Ker}{\operatorname{Ker}}
\providecommand{\cH}{\mathcal H}
\providecommand{\cO}{\mathcal O}
\providecommand{\cP}{\mathcal P}
\providecommand{\cE}{\mathcal E}
\providecommand{\cF}{\mathcal F}
\providecommand{\cG}{\mathcal G}
\providecommand{\cD}{\mathcal D}
\providecommand{\cT}{\mathcal T}
\providecommand{\cL}{\mathcal L}
\providecommand{\RHom}{\mathrm R\!\operatorname{Hom}}
\providecommand{\RGamma}{\mathrm R\Gamma}
\providecommand{\dlim}{\operatorname*{lim}}
\providecommand{\plim}{\operatorname*{lim}}

% Reviewed French continuation definitions, retained verbatim.
% Additions for Expose II, Appendix II
\providecommand{\cO}{\mathcal O}
\providecommand{\cE}{\mathscr E}
\providecommand{\cF}{\mathscr F}
\providecommand{\cG}{\mathscr G}
\providecommand{\cH}{\mathcal H}
\providecommand{\cL}{\mathscr L}
\providecommand{\cP}{\mathscr P}
\providecommand{\cM}{\mathscr M}
\providecommand{\cD}{\mathcal D}
\providecommand{\R}{\mathbb R}
\providecommand{\C}{\mathbb C}
\providecommand{\N}{\mathbb N}
\providecommand{\Zell}{\mathbb Z_\ell}
\providecommand{\Qell}{\mathbb Q_\ell}
\providecommand{\C}{\mathbb C}
\providecommand{\Ch}{\operatorname{Ch}}
\providecommand{\Sym}{\operatorname{Sym}}
\providecommand{\Td}{\operatorname{Td}}
\providecommand{\Diff}{\operatorname{Diff}}
\providecommand{\Homop}{\operatorname{Hom}}
\providecommand{\Spec}{\operatorname{Spec}}
\providecommand{\pr}{\operatorname{pr}}
\providecommand{\ind}{\operatorname{ind}}
\providecommand{\soft}{\operatorname{soft}}
\providecommand{\Ban}{\operatorname{Ban}}
\providecommand{\Lop}{\operatorname{L}}
\providecommand{\im}{\operatorname{Im}}
\providecommand{\Parf}{\operatorname{Parf}}
\renewcommand{\id}{\operatorname{id}}

\newpage

